
\paragraph{\textbf{Proposition 3.12.}}
\SGARefTarget{sga3:xiv:3.12}
{\itshape Let \(G\) be a smooth \(S\)-prescheme in groups of finite
presentation over \(S\). Let \(H\) and \(K\) be two subpreschemes
in groups, smooth and of finite presentation over \(S\), with
\(K\) having connected fibers, and let
\(\mathfrak g,\mathfrak h,\mathfrak k\) be the corresponding Lie
algebras. Suppose that one of the following conditions holds on
their geometric fibers:

\smallskip
\noindent
\textup{a)}
For every \(s\in S\), \(\mathfrak k_{\bar s}\) contains a regular
element of \(\mathfrak g_{\bar s}\), and
\(\mathfrak h_{\bar s}\) contains a Cartan subalgebra of
\(\mathfrak g_{\bar s}\).

\smallskip
\noindent
\textup{b)}
For every \(s\in S\), \(\mathfrak k_{\bar s}\) contains a Cartan
subalgebra of \(\mathfrak g_{\bar s}\).

\smallskip
\noindent
Under these conditions, \(H\supset K\) if and only if
\(\mathfrak h\supset\mathfrak k\).}

It remains only to prove that
\(\mathfrak h\supset\mathfrak k\) implies \(H\supset K\). In case
\textup{b)}, the inclusion
\(\mathfrak h\supset\mathfrak k\) places us under the hypotheses of
case \textup{a)}, so it is enough to prove \textup{a)}. Proceeding
as in \SGARefLink{sga3:xiv:3.2}{3.2}\textup{(a)}, by reduction to the case where \(S\) is
local artinian, we reduce to the case in which there exists a
section \(a\) of \(\mathfrak k\) that is quasi-regular in
\(\mathfrak g\). Arguing then as in \SGARefLink{sga3:xiii:corollary:5:5}{XIII~5.5} reduces us to the
following statement.

\paragraph{\textbf{Corollary 3.13.}}
\SGARefTarget{sga3:xiv:3.13}
{\itshape Let \(G\) be a smooth \(S\)-prescheme in groups of finite
presentation over \(S\), let \(H\subset G\) be a subprescheme in
groups smooth and of finite presentation over \(S\), and let
\(\mathfrak g\) and \(\mathfrak h\) be their Lie algebras. Let
\(a\) be a section of \(\mathfrak h\) that is quasi-regular in
\(\mathfrak g\). Suppose that for every \(s\in S\), the geometric
fiber \(\mathfrak h_{\bar s}\) contains a Cartan subalgebra of
\(\mathfrak g_{\bar s}\). Put
\[
  M_a=\operatorname{Transp}_G(a,\mathfrak h).
\]
This is a closed subprescheme of \(G\), smooth over \(S\)
\textup{(cf.\ \SGARefLink{sga3:xiv:3.1}{3.1})}. Hence \(M_a^0\), the union in the fibers of
\(M_a\) of the connected components containing the identity, is
an open subset of \(M_a\), endowed with the structure induced from
\(M_a\). Then
\[
  H^0=M_a^0.
\]}

Clearly \(H\subset M_a\), and therefore
\(H^0\subset M_a^0\). Since this is an inclusion of preschemes
smooth over \(S\), it is enough to check equality on the fibers.
This reduces us to the case where \(S\) is the spectrum of an
algebraically closed field, which was treated in \SGARefLink{sga3:xiii:corollary:5:4}{XIII~5.4}.

\paragraph{\textbf{Corollary 3.14.}}
\SGARefTarget{sga3:xiv:3.14}
{\itshape Let \(G\) be a smooth \(S\)-prescheme in groups of finite
presentation, let \(D\) be a subgroup of type \((\mathrm C)\) of
\(G\), and let \(H\subset G\) be a subprescheme in groups smooth
and of finite presentation over \(S\). If \(\mathfrak d\) and
\(\mathfrak h\) are their Lie algebras, then
\[
  \operatorname{Transp}_G(D,H)
  =
  \operatorname{Transp}_G(\mathfrak d,\mathfrak h),
\]
and this functor is representable by a closed subprescheme of
\(G\), smooth over \(S\).}

Indeed, the identity between the two transporters follows from
\SGARefLink{sga3:xiv:3.12}{3.12}\textup{(b)}, after which \SGARefLink{sga3:xiv:3.2}{3.2} applies.

\paragraph{\textbf{Corollary 3.15.}}
\SGARefTarget{sga3:xiv:3.15}
{\itshape Let \(G\) and \(H\) be as in \SGARefLink{sga3:xiv:3.14}{3.14}, and suppose that for
every \(s\in S\), the geometric fiber \(\mathfrak h_{\bar s}\)
contains a Cartan subalgebra of \(\mathfrak g_{\bar s}\). Suppose
also that \(\mathfrak g\) admits a Cartan subalgebra fpqc locally.
Then, étale locally, \(H\) contains a subgroup of type
\((\mathrm C)\) of \(G\).}

By \SGARefLink{sga3:xiv:3.7}{3.7} and \SGARefLink{sga3:xiv:3.9}{3.9}\textup{(a)}, \(G\) admits a subgroup of type
\((\mathrm C)\) étale locally. We may therefore suppose that \(G\)
admits such a subgroup \(D\). The hypothesis on \(\mathfrak h\)
then also says that the structure morphism of the transporter in
\SGARefLink{sga3:xiv:3.14}{3.14} is surjective, by the conjugacy theorem \SGARefLink{sga3:xiii:theorem:6:1}{XIII~6.1}\textup{(a)}.
The conclusion follows from Hensel's lemma \SGARefLink{sga3:xi:proposition:1:10}{XI~1.10}.
